\documentclass[12pt]{article}
\usepackage{graphicx}
\usepackage{booktabs}
\usepackage{geometry}
\geometry{margin=1in}
\title{Survival Analysis of Student Dropout: A Staged Cohort Approach}
\author{Bayes Business School}
\date{\today}

\begin{document}
\maketitle

\section{Introduction}
Student dropout is a critical challenge in higher education, with significant implications for both institutions and learners. Traditional predictive models often treat dropout as a static binary outcome, but this approach can obscure the temporal and structural nature of student disengagement. In this report, we employ a survival analysis framework, leveraging staged cohort modeling to capture the dynamics of dropout risk as students progress through key programme milestones.

\section{Rationale for Survival Analysis}
Unlike standard classification, survival analysis explicitly models the timing and stages of dropout. In our dataset, missingness is not random: students who leave the programme cease to generate new data, resulting in left-censored records. By aligning our analysis with tuition-fee checkpoints and academic milestones, we can:
\begin{itemize}
    \item Identify when students are most at risk of dropping out
    \item Tailor interventions to specific stages
    \item Avoid misleading imputations for data that could not have existed
\end{itemize}

\section{Staged Cohort Construction}
We define four longitudinal stages, each corresponding to a set of available features and a tuition-fee checkpoint:

\begin{table}[h!]
\centering
\begin{tabular}{lll}
\toprule
Stage & Columns Available & Cohort \\
\midrule
Stage 1 & Session 1, 2; Test 1; Forum/Office-Hours (amortised $\div$ 4) & Full cohort \\
Stage 2 & + Session 3, 4; Test 2; Forum/Office-Hours ($\div$ 2) & Survived Stage 1 \\
Stage 3 & + Session 5; Test 3; Forum/Office-Hours ($\times 3/4$) & Survived Stage 2 \\
Stage 4 & + Ind CW, Group CW, Final Grade; Full Forum/Office-Hours & Survived Stage 3 \\
\bottomrule
\end{tabular}
\caption{Programme stages and available features at each checkpoint.}
\end{table}

\section{Kaplan--Meier Survival Curve}
To visualize the proportion of students remaining at each stage, we construct a Kaplan--Meier style survival curve. This approach mirrors clinical trials, tracking the fraction of the cohort still enrolled after each milestone.

% Placeholder for survival curve figure
\begin{figure}[h!]
    \centering
    \includegraphics[width=0.8\textwidth]{fig_survival_curve.png}
    \caption{Kaplan--Meier style survival curve showing student retention at each programme stage.}
\end{figure}

\section{Staged Predictive Modeling}
At each stage, we fit a dedicated predictive model using only the data available up to that point:
\begin{itemize}
    \item \textbf{Stage 1:} Logistic Regression (interpretable, low complexity)
    \item \textbf{Stage 2:} Random Forest (captures non-linearities)
    \item \textbf{Stage 3:} Histogram Gradient Boosting (handles missingness natively)
    \item \textbf{Stage 4:} Stacking Ensemble (leverages all information)
\end{itemize}
Each model is trained on the sub-cohort that survived to that stage, predicting subsequent dropout.

\subsection{Cross-Stage Performance}
\begin{table}[h!]
\centering
\begin{tabular}{lcccc}
\toprule
Stage & Model & ROC-AUC (mean $\pm$ std) & PR-AUC (mean $\pm$ std) & Conditional Dropout Rate \\
\midrule
% Fill with actual results from stage_results
Stage 1 & Logistic Regression & 0.XX $\pm$ 0.XX & 0.XX $\pm$ 0.XX & XX\% \\
Stage 2 & Random Forest & 0.XX $\pm$ 0.XX & 0.XX $\pm$ 0.XX & XX\% \\
Stage 3 & Hist. Grad. Boosting & 0.XX $\pm$ 0.XX & 0.XX $\pm$ 0.XX & XX\% \\
Stage 4 & Stacking Ensemble & 0.XX $\pm$ 0.XX & 0.XX $\pm$ 0.XX & XX\% \\
\bottomrule
\end{tabular}
\caption{Model performance and dropout rates at each stage.}
\end{table}

% Placeholder for performance barplots
\begin{figure}[h!]
    \centering
    \includegraphics[width=0.8\textwidth]{fig_stage_performance.png}
    \caption{Barplots of ROC-AUC and PR-AUC for each stage.}
\end{figure}

\section{Interpretation and Discussion}
\begin{itemize}
    \item \textbf{Structural missingness:} Dropout is revealed by the absence of later-stage data, not just by the target label.
    \item \textbf{Cohort enrichment:} As students progress, the cohort shrinks but the available data per student becomes richer, justifying more complex models at later stages.
    \item \textbf{Checkpoint alignment:} Dropout spikes align with tuition-fee deadlines, making stage boundaries actionable for intervention.
    \item \textbf{Amortisation:} Forum and office-hour engagement are divided by the fraction of the programme elapsed, preventing information leakage.
    \item \textbf{Actionable metrics:} Reporting conditional dropout rates at each stage enables targeted support for at-risk students.
\end{itemize}

\section{Limitations and Future Work}
\begin{itemize}
    \item Confidence intervals for the survival curve can be added using the \texttt{lifelines} package.
    \item Cox proportional hazards models could provide interpretable continuous-time risk estimates.
    \item Cascading risk scores from early stages into later models may improve accuracy.
    \item Clustering within each stage could reveal distinct engagement personas.
\end{itemize}

\section{Conclusion}
A staged survival analysis approach provides a principled, interpretable, and actionable framework for understanding and predicting student dropout. By aligning models with programme structure and handling missingness as signal, we enable more effective early warning and intervention strategies.

\end{document}
